\section{Version history}\label{app:version-history}

Version~1 appeared on Zenodo on 20 April 2026, Version~2 on Qeios on 9 June
2026, and Version~3 on 2 October 2026. Each subsection lists the differences
from the preceding version.

\subsection{Version 3}

Version~3 keeps every numbered result of Version~2 with its number, label and
type, in the same order. The titles of Remarks~\ref{rem:repaired-alignments}
and~\ref{rem:three-repaired-alignments} (``corrected'' became ``excluded''
alignments) and of Remark~\ref{rem:recognition-open} changed.

\paragraph{Corrected statements.} Definition~\ref{def:level-trichotomy}
qualifies lossy forgetting, and Definition~\ref{def:selected-lift-atlas}
qualifies non-uniqueness and lossiness.
Definition~\ref{prop:forgetting-lifts} makes the loss-record and
uniqueness-proof requirements explicit; Version~2 already admitted
lossless forgetting and unique lifts when proved and recorded.
Proposition~\ref{prop:boundary-distinctions} is stated as the failure of
\(W_i\) for a cluster built from the bare boundary datum.
Proposition~\ref{prop:non-overbreadth} no longer claims that unaudited
material cannot enter \FATCD{}, which fails for unclaimed records.
Propositions~\ref{prop:probability-closure}--\ref{prop:agency-closure}
concluded that a feature of a recognized form classifies as a \(P_i\)
projection, which needs attachment data; they now license label (i) or (v),
or another listed label, for every well-formed decomposition record.
Definition~\ref{def:channel-status} replaces five conditions that were not
shown to be exclusive and exhaustive by an ordered rule.
Remark~\ref{rem:recognition-open} left open whether the recognition
hypothesis holds on \FATCD{}; Version~3 proves that it does not
(\(D_{\mathrm{norm}}\in\FATCD\setminus\FATCD_{\mathrm{RH}}\)), and
Theorem~\ref{thm:scoped-exact-six}, whose conclusion Version~2 named
\(\operatorname{ScopedExactSix}(\FATCD)\), now concludes
\(\operatorname{ScopedExactSix}(\FATCD_{\mathrm{RH}})\);
Definition~\ref{def:scoped-exact-six-question} records that the unrestricted
statement fails. Proposition~\ref{thm:integrated-stress} replaces the
readiness claim of Version~2 by clauses (a)--(e). The scope statements of
Remarks~\ref{rem:typed-non-collapse-vs-independence},
\ref{prop:activation-thresholds}, \ref{rem:classification-scope},
\ref{prop:dependency-stress} and~\ref{prop:residual-stress} were corrected
in the same way.

\paragraph{Added hypotheses and strengthened results.}
Definition~\ref{def:fatcd} requires realization relations to list exactly the
satisfying candidates, and Definition~\ref{def:threshold-attachment} requires
nonclaim, collapse-case and witness-schema annotations targeting the cluster.
Theorem~\ref{thm:upper-bound-classification} holds for every well-formed
decomposition record and adds a pointwise form;
Theorem~\ref{thm:six-channel-exactness} adds clause (iv), independence of the
status sequence from the decomposition record;
Proposition~\ref{prop:channel-vs-activation} realizes every active-channel
count from \(0\) to \(6\) in \(\FATCD_{\mathrm{RH}}\), and all five
statuses in one description; Proposition~\ref{prop:non-uniqueness}
exhibits distinct well-formed records in \(\FATCD_{\mathrm{RH}}\); Proposition~\ref{prop:typed-non-collapse}
covers every \(W_k\); and Theorem~\ref{thm:scoped-exact-six} adds clauses (d)
and (e).

\paragraph{Definitions and scope.} Version~3 specifies the witness
predicates through the \emph{Witness types} table, audit honesty through
Definition~\ref{def:audited}, and licensed labels through
Definition~\ref{def:residual-scheme}. It distinguishes kind residuals
from record residuals and extends recognition to bookkeeping annotations,
with alternatives (R1)--(R5). Category (xi), called a genuine seventh
irreducible typed role in Version~2, is a recorded screen, not proof
of a new role. Remark~\ref{rem:recognition-open} adds the screen example
and conditional closed-formula wrapping constructions, showing that
recognition depends on presentation.

\paragraph{Additions and removals.} Version~3 adds Subsection~\ref{subsec:objects-at-a-glance}
with its vocabulary, symbols and example index; a running example; figures of
the section dependencies and of the three levels; the label-criteria table;
and collects the Lean-counterpart notes, formerly footnotes to the
results, in Appendix~\ref{app:result-notes}. The
nonclaim ``no surrogate terminal evidence'' and the statement that
\(\mathrm{P6}_{\mathrm{drive}}\) is a face of \(\Psix{}\) were removed.
``Regressions'' and ``guardrails'' are called misreadings and checks, and
Section~\ref{sec:integrated-stress} is titled accordingly.

\subsection{Version 2}

Version~2 made the upper bound conditional on an explicit recognition
hypothesis: it added Definition~\ref{def:recognition-hypothesis} and
Remark~\ref{rem:recognition-open}, so that the upper-bound classification of
Version~1 (Theorem~6.7) became Theorem~\ref{thm:upper-bound-classification} and
its scope remark became Remark~\ref{rem:classification-scope}, and the absence
of a seventh role in the terminal theorem was stated under that hypothesis.
Statements of Section~\ref{sec:admissibility-meta-theory} that fix
admissibility conditions rather than assert consequences were stated as
definitions (Definitions~3.1, 3.5, 3.6, 3.9 and~3.11) or as a remark
(Remark~3.8), the stress checks of Section~\ref{sec:integrated-stress} as
Remarks~8.1--8.6, and the integrated stress theorem as
Proposition~\ref{thm:integrated-stress}. The proof of
Proposition~\ref{prop:boundary-distinctions} was completed over all fifteen
rows of the boundary matrix, corrections from external review were
incorporated, and notes on the Lean counterpart of each result were added,
with an appendix describing the Lean development.
